\section{Inspiración para la investigación y desafíos de despliegue}
\subsection{La fusión entre el autosupervisado y la validación en producción}
El ponente reiteró el flujo de trabajo de «Self-Feeding Chatbot»: el modelo se etiqueta a sí mismo después de la conversación, genera una verified response y luego incorpora las respuestas de alta calidad al entrenamiento, con lo que logra online self-improvement (SRT 774-884). La diferencia entre este mecanismo y el entrenamiento self-rewarding es que la señal de supervisión proviene de interacciones reales con usuarios, y no de tasks generadas de manera pura dentro del conjunto de entrenamiento.

\begin{warningbox}{Riesgo de robustez en el despliegue}
La hallucination sigue siendo una threat: si el judge considera razonable una respuesta y la hallucination resulta ser autoconsistente, todo el proceso de self-improvement puede entrar en un ciclo de amplificación de una «falacia lógica reforzada». Es necesario introducir un evaluator externo o evidencia multimodal para acotar este tipo de errores.
\end{warningbox}

\subsection{Problemas de frontera y direcciones futuras}
Weston mencionó varias direcciones futuras, entre ellas inference time compute for evaluation, interpretability y sistemas «self-aware». En concreto, dijo: «Reason by actually interacting with people in the world, internet or itself» (SRT 3282-3284), y señaló la relación entre el presupuesto de compute de la self evaluation y el reasoning understand (SRT 3236-3252). El autoconocimiento (self-aware), aunque sigue siendo algo impreciso, implica al menos que el modelo pueda buscar ayuda externa de manera activa cuando carece de un judge.

\begin{importantbox}{Tres prioridades del trabajo futuro}
\begin{itemize}
    \item El inference-time compute budget del razonamiento y la evaluación.
    \item Un pipeline de interpretation transparente, que permita a las personas hacer trace del razonamiento.
    \item Sistemas self-aware: el modelo sabe buscar un external agent cuando la uncertainty es alta.
\end{itemize}
\end{importantbox}

\subsection{Resumen de la sección}
La investigación debe equilibrar el self-training offline con el verified feedback online, y buscar un balance entre hallucination, compute e interpretability. El objetivo final es construir un modelo que razone y a la vez sea capaz de autoevaluarse.

